% ============================================================
%  MTH636 -- Game Theory | Lecture 9
%  Mixed Strategies and Expected Utility
% ============================================================
\documentclass[12pt, a4paper]{article}

% ============================================================
%  MTH636 -- Game Theory
%  Notes Preamble — Color-Coded, TikZ-Enhanced
% ============================================================

% --- Core packages ---
\usepackage[T1]{fontenc}
\usepackage[utf8]{inputenc}
\usepackage{lmodern}
\usepackage[margin=2.5cm, top=2.8cm, bottom=2.8cm]{geometry}
\usepackage{microtype}

% --- Math ---
\usepackage{amsmath, amssymb, amsthm, mathtools}
\usepackage{bm}        % bold math
\usepackage{cancel}    % strike-through in math

% --- Color & Boxes ---
\usepackage[dvipsnames, svgnames, x11names, table]{xcolor}
\usepackage[most]{tcolorbox}
\tcbuselibrary{breakable, skins, theorems}

% --- TikZ & Diagrams ---
\usepackage{tikz}
\usepackage{tikz-qtree}
\usetikzlibrary{arrows.meta, positioning, shapes.geometric, fit, backgrounds,
                decorations.pathreplacing, calc, matrix, trees}
\usepackage{pgfplots}
\pgfplotsset{compat=1.18}

% --- Tables ---
\usepackage{booktabs}
\usepackage{multirow}
\usepackage{array}
\usepackage{tabularx}
\usepackage{makecell}

% --- Typography & Layout ---
\usepackage{enumitem}
\usepackage{parskip}
\usepackage{fancyhdr}
\usepackage{titlesec}
\usepackage{titletoc}
\usepackage{hyperref}
\usepackage{graphicx}
\usepackage{float}
\usepackage{caption}

% --- Color Palette ---
\definecolor{defblue}{HTML}{1A5276}      % Definition header
\definecolor{defbluebg}{HTML}{D6EAF8}    % Definition background
\definecolor{thmgreen}{HTML}{1E8449}     % Theorem header
\definecolor{thmbg}{HTML}{D5F5E3}        % Theorem background
\definecolor{exorange}{HTML}{B7770D}     % Example header
\definecolor{exbg}{HTML}{FEF9E7}         % Example background
\definecolor{warnred}{HTML}{922B21}      % Warning/Key Insight header
\definecolor{warnbg}{HTML}{FADBD8}       % Warning background
\definecolor{noticegray}{HTML}{2C3E50}   % Note header
\definecolor{noticebg}{HTML}{EBF5FB}     % Note background
\definecolor{taketawayyellow}{HTML}{7D6608} % Key Takeaway header
\definecolor{takeawaybg}{HTML}{FDFEFE}   % Takeaway background
\definecolor{accentpurple}{HTML}{6C3483} % Accent / proof
\definecolor{lightgray}{HTML}{F2F3F4}    % shading

% ============================================================
%  TCOLORBOX ENVIRONMENTS
% ============================================================

% -- Definition (Blue) --
\newtcolorbox{defbox}[2][]{
  enhanced, breakable,
  colback=defbluebg, colframe=defblue, coltitle=white,
  fonttitle=\bfseries\sffamily,
  title={Definition: #2},
  attach boxed title to top left={yshift=-2mm, xshift=4mm},
  boxed title style={colback=defblue, rounded corners},
  arc=4pt, boxrule=1.2pt,
  left=6pt, right=6pt, top=8pt, bottom=6pt,
  #1
}

% -- Theorem / Proposition (Green) --
\newtcolorbox{thmbox}[2][]{
  enhanced, breakable,
  colback=thmbg, colframe=thmgreen, coltitle=white,
  fonttitle=\bfseries\sffamily,
  title={Theorem: #2},
  attach boxed title to top left={yshift=-2mm, xshift=4mm},
  boxed title style={colback=thmgreen, rounded corners},
  arc=4pt, boxrule=1.2pt,
  left=6pt, right=6pt, top=8pt, bottom=6pt,
  #1
}

% -- Proposition (Green variant) --
\newtcolorbox{propbox}[2][]{
  enhanced, breakable,
  colback=thmbg, colframe=thmgreen!70!black, coltitle=white,
  fonttitle=\bfseries\sffamily,
  title={Proposition: #2},
  attach boxed title to top left={yshift=-2mm, xshift=4mm},
  boxed title style={colback=thmgreen!70!black, rounded corners},
  arc=4pt, boxrule=1.2pt,
  left=6pt, right=6pt, top=8pt, bottom=6pt,
  #1
}

% -- Example (Orange/Yellow) --
\newtcolorbox{exbox}[2][]{
  enhanced, breakable,
  colback=exbg, colframe=exorange, coltitle=white,
  fonttitle=\bfseries\sffamily,
  title={Example: #2},
  attach boxed title to top left={yshift=-2mm, xshift=4mm},
  boxed title style={colback=exorange, rounded corners},
  arc=4pt, boxrule=1.2pt,
  left=6pt, right=6pt, top=8pt, bottom=6pt,
  #1
}

% -- Key Insight / Warning (Red) --
\newtcolorbox{warnbox}[2][]{
  enhanced, breakable,
  colback=warnbg, colframe=warnred, coltitle=white,
  fonttitle=\bfseries\sffamily,
  title={Key Insight: #2},
  attach boxed title to top left={yshift=-2mm, xshift=4mm},
  boxed title style={colback=warnred, rounded corners},
  arc=4pt, boxrule=1.2pt,
  left=6pt, right=6pt, top=8pt, bottom=6pt,
  #1
}

% -- Note (Gray-Blue) --
\newtcolorbox{notebox}[1][]{
  enhanced, breakable,
  colback=noticebg, colframe=noticegray!60, coltitle=white,
  fonttitle=\bfseries\sffamily,
  title={Note},
  attach boxed title to top left={yshift=-2mm, xshift=4mm},
  boxed title style={colback=noticegray!60, rounded corners},
  arc=4pt, boxrule=0.8pt,
  left=6pt, right=6pt, top=8pt, bottom=6pt,
  #1
}

% -- Key Takeaway --
\newtcolorbox{takebox}[1][]{
  enhanced, breakable,
  colback=takeawaybg, colframe=taketawayyellow, coltitle=white,
  fonttitle=\bfseries\sffamily,
  title={Key Takeaway},
  attach boxed title to top left={yshift=-2mm, xshift=4mm},
  boxed title style={colback=taketawayyellow, rounded corners},
  arc=4pt, boxrule=1.2pt,
  left=6pt, right=6pt, top=8pt, bottom=6pt,
  #1
}

% -- Intuition --
\newtcolorbox{intubox}[1][]{
  enhanced, breakable,
  colback=accentpurple!8, colframe=accentpurple!60, coltitle=white,
  fonttitle=\bfseries\sffamily,
  title={Intuition},
  attach boxed title to top left={yshift=-2mm, xshift=4mm},
  boxed title style={colback=accentpurple!80, rounded corners},
  arc=4pt, boxrule=0.8pt,
  left=6pt, right=6pt, top=8pt, bottom=6pt,
  #1
}

% ============================================================
%  PAYOFF MATRIX HELPER
%  Usage: inside tabular, use \payoff{a}{b} for (a,b)
% ============================================================
\newcommand{\payoff}[2]{$(#1,\;#2)$}
\newcommand{\nash}[2]{\underline{$#1$},\underline{$#2$}}  % underline NE payoffs
\newcommand{\Nset}{\mathbb{N}}
\newcommand{\Rset}{\mathbb{R}}
\newcommand{\Sset}{\mathcal{S}}
\newcommand{\Ui}{u_i}

% ============================================================
%  SECTION FORMATTING
% ============================================================
\titleformat{\section}
  {\large\bfseries\sffamily\color{defblue}}
  {\thesection}{1em}{}[\titlerule]

\titleformat{\subsection}
  {\normalsize\bfseries\sffamily\color{thmgreen!80!black}}
  {\thesubsection}{1em}{}

\titleformat{\subsubsection}
  {\normalsize\itshape\sffamily\color{exorange!80!black}}
  {\thesubsubsection}{1em}{}

% ============================================================
%  HEADER / FOOTER
% ============================================================
\pagestyle{fancy}
\fancyhf{}
\fancyhead[L]{\sffamily\small\color{defblue} MTH636 -- Game Theory}
\fancyhead[R]{\sffamily\small\color{defblue} \leftmark}
\fancyfoot[C]{\sffamily\small\thepage}
\renewcommand{\headrulewidth}{0.4pt}

% ============================================================
%  HYPERREF SETUP
% ============================================================
\hypersetup{
  colorlinks=true,
  linkcolor=defblue,
  urlcolor=thmgreen,
  citecolor=accentpurple,
  bookmarksnumbered=true,
  pdftitle={MTH636 Game Theory Notes},
  pdfauthor={MTH636 Student Notes}
}

% ============================================================
%  MISC HELPERS
% ============================================================
\newcommand{\BR}{\text{BR}}        % Best Response
\newcommand{\NE}{\text{NE}}        % Nash Equilibrium
\newcommand{\SPNE}{\text{SPNE}}    % Subgame Perfect NE
\newcommand{\Dom}{\text{dom}}      % Dominance
\newcommand{\mixs}[1]{\sigma_{#1}} % Mixed strategy
\newcommand{\pures}[1]{s_{#1}}     % Pure strategy


% ============================================================
\begin{document}

% Title
\begin{center}
  {\LARGE\bfseries\sffamily\color{defblue} MTH636 -- Game Theory}\\[6pt]
  {\Large\sffamily Lecture 9: Mixed Strategies}\\[4pt]
\end{center}
\vspace{0.5cm}
\hrule height 1.5pt
\vspace{0.3cm}

% ============================================================
\section{Motivation for Mixed Strategies}

In many games, if players are restricted to pure, deterministic strategies, they might not be able to guarantee a high payoff, and the game might not have a Nash Equilibrium or a Value.

\begin{exbox}{Advantage of Randomization}
  Consider the following zero-sum game for Player 1:
  \begin{center}
  \renewcommand{\arraystretch}{1.3}
  \begin{tabular}{c|cc|c}
    $P_1 \setminus P_2$ & L & R & \textbf{Min} \\
    \hline
    T & 4 & 1 & 1 \\
    B & 2 & 3 & \textbf{2} \\
    \hline
  \end{tabular}
  \end{center}
  In pure strategies, Player 1's max-min value is $\underline{v}_1 = 2$ (by playing B).
  
  What if Player 1 plays a \textbf{mixed strategy}? Suppose Player 1 plays T with probability $1/4$ and B with probability $3/4$.
  \begin{itemize}
      \item If P2 plays L, P1's expected payoff is: $\frac{1}{4}(4) + \frac{3}{4}(2) = 1 + 1.5 = 2.5$
      \item If P2 plays R, P1's expected payoff is: $\frac{1}{4}(1) + \frac{3}{4}(3) = 0.25 + 2.25 = 2.5$
  \end{itemize}
  By randomizing this way, Player 1 guarantees a payoff of $2.5$, which is strictly greater than the pure max-min value of $2$.
\end{exbox}

% ============================================================
\section{Formalizing Mixed Strategies}

\begin{defbox}{Mixed Strategy}
  Let $G = \langle N, (S_i)_{i \in N}, (u_i)_{i \in N} \rangle$ be a strategic form game in which the set of strategies $S_i$ of all players is \textbf{finite}. 
  
  A \textbf{mixed strategy} of Player $i$ is a probability distribution over their pure strategies $S_i$. 
  The set of all mixed strategies of Player $i$ is denoted by $\Sigma_i$:
  \[
    \Sigma_i = \left\{ \sigma_i : S_i \to [0,1] \;\middle|\; \sum_{s_i \in S_i} \sigma_i(s_i) = 1 \right\}
  \]
\end{defbox}

If $S_i = \{T, B\}$, a mixed strategy can be written as $\sigma_1 = (p, 1-p)$, meaning play $T$ with probability $p$ and $B$ with probability $1-p$.

\begin{defbox}{Mixed Extension of a Game}
  The \textbf{mixed extension} of the game $G$ is a new game:
  \[
    \Gamma = \langle N, (\Sigma_i)_{i \in N}, (U_i)_{i \in N} \rangle
  \]
  where $\Sigma_i$ is the set of mixed strategies for Player $i$, and $U_i$ is the \textbf{expected utility} function, defined as:
  \[
    U_i(\sigma_1, \sigma_2, \dots, \sigma_n) = \sum_{(s_1, \dots, s_n) \in S} u_i(s_1, \dots, s_n) \sigma_1(s_1) \sigma_2(s_2) \cdots \sigma_n(s_n)
  \]
\end{defbox}

% ============================================================
\section{Nash Equilibrium in Mixed Strategies}

\begin{defbox}{Mixed Strategy Nash Equilibrium}
  Every Nash Equilibrium of the mixed extension $\Gamma$ is called a \textbf{Nash Equilibrium in mixed strategies} of the original game $G$.
  
  A profile $\sigma^* = (\sigma_1^*, \dots, \sigma_n^*)$ is an NE in mixed strategies if:
  \[
    U_i(\sigma^*) \geq U_i(\sigma_i, \sigma_{-i}^*) \quad \forall i \in N, \forall \sigma_i \in \Sigma_i
  \]
\end{defbox}

\begin{thmbox}{Checking NE via Pure Deviations}
  A profile $\sigma^*$ is a mixed strategy Nash Equilibrium if and only if no player can gain by deviating to a \textbf{pure strategy}:
  \[
    U_i(\sigma^*) \geq U_i(s_i, \sigma_{-i}^*) \quad \forall i \in N, \forall s_i \in S_i
  \]
\end{thmbox}

\textbf{Proof Intuition:}
Consider any mixed deviation $\sigma_i$. The expected utility $U_i(\sigma_i, \sigma_{-i}^*)$ is just a weighted average (convex combination) of the utilities of the pure strategies:
\[
  U_i(\sigma_i, \sigma_{-i}^*) = \sum_{s_i \in S_i} \sigma_i(s_i) U_i(s_i, \sigma_{-i}^*)
\]
If $U_i(\sigma^*) \geq U_i(s_i, \sigma_{-i}^*)$ for every pure strategy $s_i$, then any weighted average of these pure strategy payoffs will also be less than or equal to $U_i(\sigma^*)$.
\[
  \sum_{s_i \in S_i} \sigma_i(s_i) U_i(s_i, \sigma_{-i}^*) \leq \sum_{s_i \in S_i} \sigma_i(s_i) U_i(\sigma^*) = U_i(\sigma^*) \sum \sigma_i(s_i) = U_i(\sigma^*)
\]

% ============================================================
\section{Linearity of Expected Utility}

Expected utility is \textbf{multilinear} in the probabilities. If we hold all other players' mixed strategies fixed, the expected utility $U_i$ is a linear function of Player $j$'s mixed strategy $\sigma_j$:
\[
  U_i(\dots, \lambda\sigma_j + (1-\lambda)\sigma_j', \dots) = \lambda U_i(\dots, \sigma_j, \dots) + (1-\lambda) U_i(\dots, \sigma_j', \dots)
\]
This linearity is the fundamental property that allows us to solve for mixed equilibria using indifference conditions (as we will see).

\end{document}